% document formatting
\documentclass[10pt]{article}
\usepackage[utf8]{inputenc}
\usepackage[left=1in,right=1in,top=1in,bottom=1in]{geometry}
\usepackage[T1]{fontenc}
\usepackage{xcolor}
\usepackage[table,xcdraw]{xcolor}

% math symbols, etc.
\usepackage{amsmath, amsfonts, amssymb, amsthm}

% lists
\usepackage{enumerate}
\usepackage{tabularx}
\usepackage{multirow}

% images
\usepackage{graphicx} % for images
\usepackage{tikz}

% code blocks
\usepackage{minted, listings} 

% verbatim greek
\usepackage{alphabeta}

\newcommand{\dd}{\text{d}}

\graphicspath{{./assets/images/Week 9}}

\title{03-621 Week 9 \\ \large{Advanced Quantative Genetics}}
 
\author{Aidan Jan}

\date{\today}

\begin{document}
\maketitle

\section*{$h^2$ from Parent-Offspring Regression}
\begin{center}
    \includegraphics*[width=0.8\textwidth]{W9_1.png}
\end{center}
\begin{align*}
    \beta &= \frac{\text{Cov}(x, y)}{\text{Var}(x)} \\
    &= \frac{\sum (x_i - \bar{x})(y_i - \bar{y}) / (N - 1)}{\sum (x_i - \bar{x})^2 / (N - 1)} \\
    &= \frac{(1/2) V_A}{(1/2) V_p} \\
    &= \frac{V_A}{V_p} \\
    &= h^2
\end{align*}

\subsection*{Possible values of $h^2$ and their interpretation}
\begin{center} 
	\includegraphics*[width=0.8\textwidth]{W9_2.png} 
\end{center}
\begin{itemize}
	\item $h^2 = 1.0$: Genes with \textbf{additive} effects determine \textbf{all} of the differences in phenotypic value
	\item $1.0 > h^2 > 0$: Genes with additive effects, dominance, and/or epistasis may contribute, as well as the environment
	\item $h^2 = 0.0$: \textbf{No additive genetic contribution} to the differences in phenotypic value
\end{itemize}

\subsection*{$h^2$ for Some Human Disease Traits}
\begin{center}
    \begin{tabular}{|c|c|}
        \hline
        \textbf{Trait} & $h^2$ \\ \hline
        Cleft lip & 65 \\
        Spina bifida & 65 \\
        Club foot & 80 \\
        Pyloric stenosis & 85 \\
        Dislocation of hip & 95 \\
        Hydrocephalus & 40 \\
        Bipolar affective disorder & $\sim$100 \\
        Multiple sclerosis & $\sim$100 \\
        Schizophrenia & $\sim$100 \\
        Diabetes & 80 \\
        Epilepsy & 40 \\ \hline
    \end{tabular}
\end{center}
These particular examples are all \textbf{threshold traits:} the heritability is for the quantitative phenotype, which is the \textbf{midparent value for the underlying liability (risk)}.\\\\
\textbf{Note} that $h^2$ for a given trait can vary:
\begin{itemize}
	\item between different populations (different gene pools)
	\item between different locations (different environments)
	\item between different generations of the same population
	\begin{itemize}
	    \item Before / After selection
	    \item With changes in environment (e.g., diet)
    \end{itemize}
\end{itemize}
For the same trait and environment:
\begin{itemize}
	\item $h^2$ will be higher in a population with high genetic diversity
	\item $h^2$ will be lower in a population with low genetic diversity
\end{itemize}

\subsection*{Selection}
$h^2$ decreases when strong selection pressure reduces genetic diversity.
\begin{itemize}
	\item Consider a trait with phenotype values 4 to 16, with the (normal distribution) centered at 10.
	\item If suddenly selection pressure favors only those with above phenotype value 12, then the genotypes contributing to the next generation would be those with high phenotype value
	\item This reduces the gene pool and therefore $h^2$ decreases
\end{itemize}
Can we predict \textit{how} much the average phenotype will actually chnage under selection? \\
Can we predict the \textit{average} phenotype of progeny for any two parents?
\begin{itemize}
	\item Yes and Yes!
\end{itemize}
If we know the $h^2$ for a trait in a population, the equation for parent-offspring regression allows us to find the \textbf{average} phenotype of progeny:
\[y = h^2 x + b\]
\begin{itemize}
	\item $y$ = expected offspring value
	\item $x$ = mid-parental value
\end{itemize}
We can also obtain $h^2$ from (and apply it to)
\begin{itemize}
	\item Group Selection Experiments
	\item Threshold Traits
\end{itemize}
To do this we modify the regression equation to obtain the ``Breeder's Equation'

\subsection*{Linking $h^2$ and Selection: The Breeder's Equation}
\begin{itemize}
	\item Parent-offspring regression: $y = h^2 x + b$
	\item Rewrite in terms of differences from the original population $M$: $(y - M) = h^2 (x - M)$
	\item Average this over the selected parents and their offspring: $(M' - M) = h^2(M^* - M)$
	\begin{itemize}
	    \item $M'$ = mean of progeny
	    \item $M^*$ = mean of selected parents
    \end{itemize}
    \item \textbf{The Breeder's Equation:} $R = h^2 S$
    \begin{itemize}
	    \item $R = (M' - M)$.  ``\textbf{R}esponse to selection'' 
	    \item $S = (M^* - M)$.  ``\textbf{S}election differential''
    \end{itemize}
    \item The amount of evolutionary change ($R$) is determined by the narrow-sense heritability ($h^2$) and the strength of selection ($S$).
    \item $h^2$ can be \textbf{obtained} from selective breeding experiments \textit{or} from the incidence of a threshold trait among children of affected parents in a population
    \[h^2 = \frac{R}{S} = \frac{M' - M}{M^* - M}\]
    \item $h^2$ can be used to \textbf{predict} the phenotypic mean of a progeny population obtained by selective breeding (or natural selection)
    \[M' = M + h^2(M^* - M)\]
\end{itemize}

\subsubsection*{Example: Effect of Selection on Size of Tobacco Flowers}
We can obtain $h^2$ from selection experiments:
\begin{center} 
	\includegraphics*[width=0.8\textwidth]{W9_3.png} 
\end{center}
And we can predict the effect of selection:
\begin{center} 
	\includegraphics*[width=0.8\textwidth]{W9_4.png} 
\end{center}

\subsubsection*{For how long?  Response over multiple generations}
\begin{itemize}
	\item \textbf{strictly speaking}, an initial $h^2$ from the breeders' equation only holds for predicting a \textbf{single} generation of response.
	\begin{itemize}
	    \item Reason: allele and genotype frequencies change after selection so that $\sigma_A^2, \sigma_D^2, \sigma_I^2$, and $\sigma_{gxe}^2$ can change.  Environment can also change.
    \end{itemize}
    \item In practice, an initial $h^2$ usually does well for 5-10 generations, sometimes much longer
\end{itemize}
This depends on:
\begin{enumerate}
	\item Degree of environmental change between generations
	\item Degree of genetic change between the generation in which $h^2$ was estimated and the generation in which selection is applied.  Strong/Fast vs. Moderate/Gradual selection.
\end{enumerate}

\subsection*{$h^2$ does not always decrease after selection}
\textbf{Example: Oil content in corn}
\begin{itemize}
	\item Original $h^2$ for oil content in corn has held since 1896, or 125 generations(!)
	\item Implies gradual selection has not exhausted genetic diversity
\end{itemize}
\textbf{Example: Lactation milk yield in dairy cows}
\begin{itemize}
	\item $h^2$ for lactation milk yield in dairy cattle has actually \textit{increased} over time:
	\begin{itemize}
	    \item $\sim 25\%$ in the 1970s
	    \item $\sim 40\%$ currently
    \end{itemize}
    \item The reason for this is due to improvements in feeding and veterinary care.
    \item Leads to reduced $V_E$
\end{itemize}

\subsection*{Obtaining $h^2$ values for human traits}
\begin{enumerate}
	\item Continuous and Meristic traits
	\begin{itemize}
	    \item Phenotypic values can be measured or counted directly
	    \begin{itemize}
	        \item Use parent-offspring regression
        \end{itemize}
    \end{itemize}
    \item Threshold traits
    \begin{itemize}
	    \item We can't measure the underlying liability directly.
	    \item We can't do selection experiments
	    \item However, we can use the Breeders' Equation, adapted to incidence distributions
    \end{itemize}
\end{enumerate}

\subsection*{The Breeder's Equation for Selection}
Our examples so far have represented \textbf{Truncation Selection}
\begin{itemize}
	\item e.g., only individuals \textbf{above} a certain phenotypic value ($T$) contribute to the next generation.
\end{itemize}
Expression of Threshold Traits is equivalent to truncation selection!
\begin{itemize}
	\item The quantitative phenotype with a continuous distribution is the underlying \textbf{liability}.  The trait is expressed if liability \textbf{exceeds} a threshold.
	\item This is the same situation as threshold selection.  (Members below the threshold do not reproduce.)
	\item This equivalence gives us a way of estimating $h^2$ for the liability of threshold traits using the Breeder's Equation
\end{itemize}

\subsection*{What Proportion of the \textit{Liability} is Inherited?}
Recall:
\[h^2 = \frac{M' - M}{M^* - M} = \frac{R}{S}\]
\begin{itemize}
	\item $M$ = mean liability of the base population
	\item $M^*$ = mean liability of affected individuals
	\item $M'$ = mean liability of progeny from affected individuals
\end{itemize}
\textbf{How do we obtain $M$, $M^*$, and $M'$?}
\begin{itemize}
	\item We don't need them: only the differences $(R, S)$, which can be calculated from the proportions of affected individuals, using properties of the \textit{standard normal distribution}.
\end{itemize}
\begin{center} 
	\includegraphics*[width=0.5\textwidth]{W9_5.png} 
\end{center}
This determines what proportion of the progeny of affected individuals in the population will be above threshold and develop the trait.
\begin{itemize}
	\item Note: A \textbf{standard normal distribution} is one with mean $\mu = 0$ and variance $\sigma^2 = 1$.
	\item Any arbitrary normal distribution can be converted into a ``standard'' normal distribution by expressing the ordinate values in terms of standard deviations from the mean:
	\[z = (x - \mu) / \sigma\]
\end{itemize}

\subsection*{Standardizing the Breeder's Equation}
We express $R = M' - M$ and $S = M^* - M$ in terms of standard deviations of the mean.
\begin{itemize}
	\item Let $i = (M^* - M) / \sigma$, so $S = i\sigma$.  ($i$ stands for ``selection intensity'')
	\item Let $x_1 = (Threshold - M) / \sigma$ and $x_2 = (Threshold - M') / \sigma'$, so $R = (x_1 - x_2)\sigma$.  ($\sigma' \sim \sigma$)
\end{itemize}
Now, we get $h^2 = \frac{R}{S} = \frac{(x_1 - x_2) \sigma}{i \sigma}$
\[h^2 = \frac{x_1 - x_2}{i}\]
Now we're in business, because for a standard normal distribution,
\[P(x) = \sqrt{\frac{1}{2\pi}} \exp \left(-\frac{x^2}{2}\right)\]
The area under the curve from $-\infty$ to $X$ is
\[\Phi(X) = \int_{-\infty}^X P(x) \:\dd x\]
or, the ``cumulative distribution function''.  Similarly, the area unde the curve from $X$ to $\infty$ would be $1 - \Phi(x)$.
\begin{itemize}
	\item The \textbf{mean value} for individuals within area $1 - \Phi(X)$ is
	\[\mu^* = \frac{\exp\left(-\frac{x^2}{2}\right)}{\sqrt{2\pi} (1 - \Phi(x))}\]
\end{itemize}
Now we have what we need to obtain $h^2$ for threshold traits:
\[h^2 = \frac{x_1 - x_2}{i}\]
\begin{enumerate}
	\item If $x$ is the distance between Threshold and Mean:
	\begin{itemize}
	    \item $1 - \Phi(x)$ is the proportion of the population that exceeds the threshold and is \textbf{affected}
    \end{itemize}
    \item Given $1 - \Phi(x_1)$ and $1 - \Phi(x_2)$, we can obtain $x_1$ and $x_2$ from the distribution function $\Phi(x)$
    \item Then, we can calculate $i$:
    \[i = \mu^* = \frac{\exp \left(-\frac{x_1^2}{2}\right)}{\sqrt{2\pi} (1 - \Phi(x_1))}\]
\end{enumerate}
\textbf{Note:} in practice, if we want to find $\Phi(x)$, we would use a statistical table, or use a calculator.

\subsection*{$h^2$ for Threshold Traits \dots again}
\[h^2 = \frac{x_1 - x_2}{i}\]
Re-emphasize that:
\begin{itemize}
	\item $x_1$ is the number of standard deviations by which the Threshold exceeds the mean liability in the \textbf{base (i.e., ``general'') population}.
	\item $x_2$ is the number of standard deviations by which the Threshold exceeds the mean liability among \textbf{offspring of affected} individuals.
	\item $i$ is the number of standard deviations by which $\mu^*$ exceeds the mean liability in the \textbf{base population}.
\end{itemize}

\subsection*{Modification for Estimation for \textbf{Rare}} Threshold Traits
The $h^2$ equation above holds for \textit{common} traits where $x_2$ is calculated from progeny of \textbf{two affected parents} (so midparent liability is $\mu^* = i$)
\begin{itemize}
	\item \textbf{If the trait is rare}, only 1 parent is likely affected and has average liability $\mu^* = i$
	\item The other parent comes from the general population and has average liability $0$.
	\item Therefore, \textbf{midparent liability} is $(i + 0) / 2 = i / 2$.
\end{itemize}
\[\therefore h^2 = \frac{x_1 - x_2}{i / 2} = \frac{2(x_1 - x_2)}{i}\]

\subsubsection*{Example: Calculating $h^2$ for Threshold Traits}
For this example, consider Rheumatic Fever.  The frequency in the general population is $0.005 = 1 - \Phi(x_1)$.
\begin{itemize}
	\item So, $\Phi(x_1) = 0.995$
	\item From the distribution function we find: when $\Phi(x_1) = 0.995 \rightarrow x_1 = 2.57$.  So, $i = 2.8988$.
\end{itemize}
Frequency among offspring of affeted individuals is $0.05 = 1 - \Phi(x_2)$.
\begin{itemize}
	\item So, $\Phi(x_2) = 0.95$
	\item From the distribution function we find: when $\Phi(x_2) = 0.95 \rightarrow x_2 = 1.64$.
\end{itemize}
Therefore,
\[h^2 = \frac{x_1 - x_2}{i / 2} = \frac{2.57 - 1.64}{2.8988 / 2} = 0.64\]
So, 64\% of the liability is heritable.
\begin{itemize}
	\item Since $H^2 \geq h^2$, \textbf{at least} 64\% of the variance in liability is due to genotype.
\end{itemize}

\subsection*{Remember: Broad- versus Narrow-Sense Heritability}
\[H^2 = \frac{\sigma_g^2}{\sigma_p^2}\]
Includes:
\begin{itemize}
	\item \textbf{Additive} effects of alleles ($\sigma_A^2$)
	\item \textbf{Dominance} effects ($\sigma_D^2$)
    \item \textbf{Epistasis} effects ($\sigma_I^2$)
    \item $\sigma_g^2 = \sigma_A^2 + \sigma_D^2 + \sigma_I^2$ 
\end{itemize}
\[h^2 = \frac{\sigma_A^2}{\sigma_p^2}\]
Includes: 
\begin{itemize}
	\item Only the \textbf{additive} effects of alleles
	\item E.g., only the \textbf{additive} effects lead to heritable differences in phenotypic mean
\end{itemize}
$h^2$ is the proportion of phenotypic variance that can be used to predict the changes in the phenotypic mean after \textbf{selection on a population} or among \textbf{progeny of individuals.}
\begin{itemize}
	\item Generally, $h^2 < H^2$.
	\item $h^2 = H^2$ when \textbf{all} alleles affecting a trait are \textbf{additive} in their effects (there is no dominance or epistasis)
\end{itemize}

\subsection*{Review: Variance and Gene Number}
\[n = \frac{D^2}{8V_g}\]
The following phenomena increase when $V_g$ is increased, leading to underestimates of gene number ($n$):
\begin{itemize}
	\item Dominance
	\item Epistasis (examples: masking; multiplicative)
	\item Unequal gene/allele contribution
	\item Linkage
\end{itemize}

\subsection*{Effect of Unequal Contributions or Partially Linked Genes}
\begin{center} 
	\includegraphics*[width=0.6\textwidth]{W9_6.png} 
\end{center}
Partial linkage has similar effects because recombination in F1 germlines generates recombinant chromosomes with different contributions.


\end{document}